\documentclass[12pt,a4paper]{article}
\usepackage{../../../myconfig}

\begin{document}

% =========================================================================
% BÀI TẬP RÈN LUYỆN OXYZ (20 CÂU TỔNG HỢP)
% =========================================================================
\titlebox{Chuyên đề: Oxyz}{Bổ sung kiến thức Oxyz}{Oxyz: Bổ sung kiến thức}

\vspace{-0.25cm}

% Câu 1 (Câu 437 trong tài liệu)
% ĐÁP ÁN: A (cos phi = |2*1 + (-2)*(-1) + (-1)*1| / (sqrt(6) * sqrt(6)) = 3/6 = 1/2 => 60 độ)
\questionLinesManual{9}{1}{%
Trong không gian với hệ tọa độ $Oxyz$, cho hai mặt phẳng $(P)\colon x - 2y - z + 2 = 0$ và $(Q)\colon 2x - y + z + 1 = 0$. Góc giữa hai mặt phẳng $(P)$ và $(Q)$ bằng
\begin{tasks}(4)
    \task $60^\circ$.
    \task $90^\circ$.
    \task $30^\circ$.
    \task $120^\circ$.
\end{tasks}
}{}

\vspace{0.15cm}

% Câu 2 (Câu 440 trong tài liệu)
% ĐÁP ÁN: C (cos phi = |1*(-1) + (-1)*1 + 2*1| / (sqrt(6) * sqrt(3)) = 0... u1.u2 = -1-1+2 = 0 => 90 độ, đáp án đúng là D)
\questionLinesManual{9}{2}{%
Trong không gian với hệ tọa độ $Oxyz$, tính số đo góc giữa hai đường thẳng $d_1\colon \dfrac{x}{1} = \dfrac{y + 1}{-1} = \dfrac{z - 1}{2}$ và $d_2\colon \dfrac{x + 1}{-1} = \dfrac{y}{1} = \dfrac{z - 3}{1}$.
\begin{tasks}(4)
    \task $45^\circ$.
    \task $30^\circ$.
    \task $60^\circ$.
    \task $90^\circ$.
\end{tasks}
}{}

\vspace{0.15cm}

% Câu 3 (Câu 442 trong tài liệu)
% ĐÁP ÁN: C (u = (-1; 0; 1), n = (0; 1; -1) => sin phi = |-1| / (sqrt(2)*sqrt(2)) = 1/2 => 30 độ)
\questionLinesManual{11}{3}{%
Trong không gian với hệ tọa độ $Oxyz$, tính số đo góc giữa đường thẳng $d\colon \begin{cases} x = 2 - t \\ y = 5 \\ z = 1 + t \end{cases}$ và mặt phẳng $(P)\colon y - z + 2 = 0$.
\begin{tasks}(4)
    \task $90^\circ$.
    \task $60^\circ$.
    \task $30^\circ$.
    \task $45^\circ$.
\end{tasks}
}{}

\vspace{0.15cm}

% Câu 4 (Câu 439 trong tài liệu)
% ĐÁP ÁN: C (ud = [nP, nQ] = (0; 3; 3) cùng phương (0; 1; 1), u_Oz = (0; 0; 1) => cos = 1/sqrt(2) => 45 độ)
\questionLinesManual{11}{4}{%
Trong không gian với hệ tọa độ $Oxyz$, gọi $d$ là giao tuyến của hai mặt phẳng $(P)\colon 2x - y + z + 2017 = 0$ và $(Q)\colon x + y - z + 5 = 0$. Tính số đo góc giữa đường thẳng $d$ và trục $Oz$.
\begin{tasks}(4)
    \task $60^\circ$.
    \task $0^\circ$.
    \task $45^\circ$.
    \task $30^\circ$.
\end{tasks}
}{}

\vspace{0.15cm}

% Câu 5 (Câu 391 trong tài liệu)
% ĐÁP ÁN: A (d = |1 + 2(-2) - 2(-3) + 3| / 3 = |1 - 4 + 6 + 3| / 3 = 6/3 = 2)
\questionLinesManual{11}{5}{%
Trong không gian với hệ tọa độ $Oxyz$, cho mặt phẳng $(P)\colon x + 2y - 2z + 3 = 0$. Khoảng cách từ điểm $A(1; -2; -3)$ đến mặt phẳng $(P)$ bằng
\begin{tasks}(4)
    \task $2$.
    \task $\dfrac{2}{3}$.
    \task $\dfrac{1}{3}$.
    \task $1$.
\end{tasks}
}{}

\vspace{0.15cm}

% Câu 6 (Câu 397 trong tài liệu)
% ĐÁP ÁN: A (d = |-3| / sqrt(4 + 4 + 1) = 3/3 = 1)
\questionLinesManual{11}{6}{%
Trong không gian với hệ tọa độ $Oxyz$, tính khoảng cách từ gốc tọa độ $O(0;0;0)$ đến mặt phẳng $(P)\colon 2x + 2y + z - 3 = 0$.
\begin{tasks}(4)
    \task $1$.
    \task $\dfrac{1}{3}$.
    \task $2$.
    \task $3$.
\end{tasks}
}{}

\vspace{0.15cm}

% Câu 7 (Câu 449 trong tài liệu)
% ĐÁP ÁN: C (A(1; 0; 2) in d, u = (1; 2; 1) => AM = (1; 0; -1) => [AM, u] = (2; -2; 2) => d = sqrt(12)/sqrt(6) = sqrt(2))
\questionLinesManual{11}{7}{%
Trong không gian với hệ tọa độ $Oxyz$, khoảng cách từ điểm $M(2; 0; 1)$ đến đường thẳng $d\colon \dfrac{x - 1}{1} = \dfrac{y}{2} = \dfrac{z - 2}{1}$ bằng
\begin{tasks}(4)
    \task $\sqrt{12}$.
    \task $\sqrt{3}$.
    \task $\sqrt{2}$.
    \task $\dfrac{12}{\sqrt{6}}$.
\end{tasks}
}{}

\vspace{0.15cm}

% Câu 8 (Câu 438 trong tài liệu)
% ĐÁP ÁN: B (d = |3 - (-1)| / sqrt(1 + 4 + 4) = 4/3)
\questionLinesManual{11}{8}{%
Trong không gian với hệ tọa độ $Oxyz$, cho hai mặt phẳng song song $(P)\colon x + 2y - 2z + 3 = 0$ và $(Q)\colon x + 2y - 2z - 1 = 0$. Khoảng cách giữa hai mặt phẳng đã cho bằng
\begin{tasks}(4)
    \task $\dfrac{4}{9}$.
    \task $\dfrac{4}{3}$.
    \task $\dfrac{2}{3}$.
    \task $4$.
\end{tasks}
}{}

\vspace{0.15cm}

% Câu 9 (Câu 455 trong tài liệu)
% ĐÁP ÁN: A (u.(nP) = 2(2) + 1(-2) + 2(-1) = 0 => Delta // (P). Lay M(1; -2; 1) in Delta => d = |2(1) - 2(-2) - 1 + 1| / 3 = 6/3... doi chieu cong thuc: |2*1 - 2*(-2) - 1 + 1|/3 = 6/3 = 2. Dap an đúng: D)
\questionLinesManual{11}{9}{%
Trong không gian với hệ tọa độ $Oxyz$, cho mặt phẳng $(P)\colon 2x - 2y - z + 1 = 0$ và đường thẳng $\Delta\colon \dfrac{x - 1}{2} = \dfrac{y + 2}{1} = \dfrac{z - 1}{2}$. Khoảng cách giữa đường thẳng $\Delta$ và mặt phẳng $(P)$ bằng
\begin{tasks}(4)
    \task $\dfrac{1}{3}$.
    \task $\dfrac{5}{3}$.
    \task $\dfrac{2}{3}$.
    \task $2$.
\end{tasks}
}{}

\vspace{0.15cm}

% Câu 10 (Câu 456 trong tài liệu)
% ĐÁP ÁN: B (M(3; -2; -1), N(0; 1; 2) => MN = (-3; 3; 3); u1 = (-4; 1; 1), u2 = (-6; 1; 2) => [u1, u2] = (1; 2; 2) => d = |-3 + 6 + 6| / 3 = 9/3 = 3)
\questionLinesManual{11}{10}{%
Trong không gian với hệ tọa độ $Oxyz$, cho hai đường thẳng chéo nhau $\Delta_1\colon \dfrac{x - 3}{-4} = \dfrac{y + 2}{1} = \dfrac{z + 1}{1}$ và $\Delta_2\colon \dfrac{x}{-6} = \dfrac{y - 1}{1} = \dfrac{z - 2}{2}$. Khoảng cách giữa hai đường thẳng $\Delta_1$ và $\Delta_2$ bằng
\begin{tasks}(4)
    \task $\dfrac{27}{\sqrt{209}}$.
    \task $3$.
    \task $1$.
    \task $\dfrac{5}{3}$.
\end{tasks}
}{}

\vspace{0.15cm}

% Câu 11 (Câu 448 trong tài liệu)
% ĐÁP ÁN: D (R = d(I, P) = |12(4) - 5(-2) - 19| / sqrt(12^2 + (-5)^2) = |48 + 10 - 19| / 13 = 39/13 = 3)
\questionLinesManual{11}{11}{%
Trong không gian với hệ tọa độ $Oxyz$, cho mặt cầu $(S)$ có tâm $I(4; 2; -2)$ tiếp xúc với mặt phẳng $(P)\colon 12x - 5z - 19 = 0$. Khi đó bán kính $R$ của mặt cầu $(S)$ bằng
\begin{tasks}(4)
    \task $39$.
    \task $\dfrac{39}{\sqrt{13}}$.
    \task $13$.
    \task $3$.
\end{tasks}
}{}

\vspace{0.15cm}

% Câu 12 (Câu 445 trong tài liệu)
% ĐÁP ÁN: B (M in tia Oz => M(0; 0; c), c >= 0. d(M, P) = |c + 6| / 3 = 3 => |c + 6| = 9 => c = 3 (nhan), c = -15 (loai) => M(0; 0; 3))
\questionLinesManual{11}{12}{%
Trong không gian với hệ tọa độ $Oxyz$, cho mặt phẳng $(P)\colon 2x + 2y + z + 6 = 0$. Tìm tọa độ điểm $M$ thuộc tia $Oz$ sao cho khoảng cách từ $M$ đến mặt phẳng $(P)$ bằng $3$.
\begin{tasks}(2)
    \task $M(0; 0; 21)$.
    \task $M(0; 0; 3)$.
    \task $M(0; 0; 3)$ hoặc $M(0; 0; -15)$.
    \task $M(0; 0; -15)$.
\end{tasks}
}{}

\vspace{0.15cm}

% Câu 13 (Câu 469 trong tài liệu)
% ĐÁP ÁN: B (Thay A(0; -1; 2): 3(0) - (-1) + 2 - 2 = 1 != 0 => A khong thuoc; thay B(4; 1; -1): 3(4) - 1 + (-1) - 2 = 8 != 0. Tich f(A)*f(B) = 1*8 > 0 => Cung phia. Dap an dung: C)
\questionLinesManual{11}{13}{%
Trong không gian với hệ tọa độ $Oxyz$, cho hai điểm $A(0; -1; 2)$, $B(4; 1; -1)$ và mặt phẳng $(\alpha)\colon 3x - y + z - 2 = 0$. Khẳng định nào sau đây là đúng?
\begin{tasks}(2)
    \task $A \notin (\alpha)$, $B \in (\alpha)$.
    \task $A \in (\alpha)$, $B \notin (\alpha)$.
    \task $A$ và $B$ nằm về cùng một phía đối với $(\alpha)$.
    \task $A$ và $B$ nằm về hai phía khác nhau đối với $(\alpha)$.
\end{tasks}
}{}

\vspace{0.15cm}

% Câu 14 (Câu 382 trong tài liệu)
% ĐÁP ÁN: B (1/2 = -3/(-6) = 2/m != -3/(-m) => 2/m = 1/2 => m = 4. Khi m = 4 thi -3/(-4) = 3/4 != 1/2 (thoa man //))
\questionLinesManual{11}{14}{%
Trong không gian với hệ tọa độ $Oxyz$, cho mặt phẳng $(P)\colon x - 3y + 2z - 3 = 0$ và $(Q)\colon 2x - 6y + mz - m = 0$ ($m$ là tham số thực). Tìm $m$ để $(P)$ song song với $(Q)$.
\begin{tasks}(4)
    \task $m = 2$.
    \task $m = 4$.
    \task $m = -6$.
    \task $m = -10$.
\end{tasks}
}{}

\vspace{0.15cm}

% Câu 15 (Câu 403 trong tài liệu)
% ĐÁP ÁN: C (nP.nQ = 2(4) + a(-1) + 3(-(a+4)) = 0 <=> 8 - a - 3a - 12 = 0 <=> -4a - 4 = 0 <=> a = -1)
\questionLinesManual{11}{15}{%
Trong không gian với hệ tọa độ $Oxyz$, cho hai mặt phẳng $(P)\colon 2x + ay + 3z - 5 = 0$ và $(Q)\colon 4x - y - (a + 4)z + 1 = 0$. Tìm giá trị thực của tham số $a$ để $(P)$ vuông góc với $(Q)$.
\begin{tasks}(4)
    \task $a = 1$.
    \task $a = 0$.
    \task $a = -1$.
    \task $a = \dfrac{1}{3}$.
\end{tasks}
}{}

\vspace{0.15cm}

% Câu 16 (Câu 385 trong tài liệu)
% ĐÁP ÁN: A (Thay x, y, z vao P: 2(-3+t) + (2-2t) + 3(1) + 1 = -6 + 2t + 2 - 2t + 3 + 1 = 0 (dung voi moi t) => d nam trong (P))
\questionLinesManual{11}{16}{%
Trong không gian với hệ tọa độ $Oxyz$, cho mặt phẳng $(P)\colon 2x + y + 3z + 1 = 0$ và đường thẳng $d\colon \begin{cases} x = -3 + t \\ y = 2 - 2t \\ z = 1 \end{cases}$. Khẳng định nào sau đây là đúng?
\begin{tasks}(4)
    \task $d \subset (P)$.
    \task $d \perp (P)$.
    \task $d$ cắt $(P)$.
    \task $d // (P)$.
\end{tasks}
}{}

\vspace{0.15cm}

% Câu 17 (Câu 384 trong tài liệu)
% ĐÁP ÁN: B (Giai he: -3+2t = 5+t' va -2+3t = -1-4t' => t = 3, t' = -2. Thay t=3 vao d: x=3, y=7, z=18 => (3; 7; 18))
\questionLinesManual{9}{17}{%
Trong không gian với hệ tọa độ $Oxyz$, tọa độ giao điểm của hai đường thẳng cắt nhau $d\colon \begin{cases} x = -3 + 2t \\ y = -2 + 3t \\ z = 6 + 4t \end{cases}$ và $d'\colon \begin{cases} x = 5 + t' \\ y = -1 - 4t' \\ z = 20 + t' \end{cases}$ là
\begin{tasks}(4)
    \task $(5; -1; 20)$.
    \task $(3; 7; 18)$.
    \task $(-3; -2; 6)$.
    \task $(3; -2; 1)$.
\end{tasks}
}{}

\vspace{0.15cm}

% Câu 18 (Câu 409 trong tài liệu)
% ĐÁP ÁN: A (u1 = (2; 1; 2), u2 = (2; 3; 1) khong cung phuong. Giai he toa do co nghiem duy nhat t = 0, suy ra cat nhau tai (3; 0; 3))
\questionLinesManual{9}{18}{%
Trong không gian với hệ tọa độ $Oxyz$, cho hai đường thẳng $d_1\colon \dfrac{x - 1}{2} = \dfrac{y + 1}{1} = \dfrac{z - 1}{2}$ và $d_2\colon \begin{cases} x = 3 + 2t \\ y = 3t \\ z = 3 + t \end{cases}$. Vị trí tương đối giữa $d_1$ và $d_2$ là
\begin{tasks}(2)
    \task $d_1$ cắt $d_2$.
    \task $d_1 \equiv d_2$.
    \task $d_1$ và $d_2$ chéo nhau.
    \task $d_1 // d_2$.
\end{tasks}
}{}

\vspace{0.15cm}

% Câu 19 (Câu 424 trong tài liệu)
% ĐÁP ÁN: A (u = (-3; 1; -2), u' = (6; -2; 4) cung phuong (ti le -1/2). Lay M(2; -2; -1) thuoc d thay vao d' thay khong thuoc => d // d')
\questionLinesManual{11}{19}{%
Trong không gian với hệ tọa độ $Oxyz$, xét vị trí tương đối giữa hai đường thẳng $d\colon \dfrac{x - 2}{-3} = \dfrac{y + 2}{1} = \dfrac{z + 1}{-2}$ và $d'\colon \dfrac{x}{6} = \dfrac{y - 4}{-2} = \dfrac{z - 2}{4}$. Mệnh đề nào sau đây là đúng?
\begin{tasks}(4)
    \task $d // d'$.
    \task $d \equiv d'$.
    \task $d$ cắt $d'$.
    \task $d$ và $d'$ chéo nhau.
\end{tasks}
}{}

\vspace{0.15cm}

% Câu 20 (Câu 470 trong tài liệu)
% ĐÁP ÁN: B (I1(2; 0; -1), R1 = 4; I2(-3; 2; 0), R2 = 1. I1I2 = sqrt(25 + 4 + 1) = sqrt(30) approx 5.477. R1 + R2 = 5 < sqrt(30) => O ngoai nhau, khong co diem chung)
\questionLinesManual{11}{20}{%
Trong không gian với hệ tọa độ $Oxyz$, cho hai mặt cầu $(S_1)\colon (x - 2)^2 + y^2 + (z + 1)^2 = 16$ và $(S_2)\colon (x + 3)^2 + (y - 2)^2 + z^2 = 1$. Khẳng định nào sau đây là đúng?
\begin{tasks}(2)
    \task $(S_1)$ và $(S_2)$ cắt nhau.
    \task $(S_1)$ và $(S_2)$ không có điểm chung.
    \task $(S_1)$ và $(S_2)$ tiếp xúc trong.
    \task $(S_1)$ và $(S_2)$ tiếp xúc ngoài.
\end{tasks}
}{}

\drawdashlinesfull

% 1A-2D-3C-4C-5A-6A-7C-8B-9D-10B-11D-12B-13C-14B-15C-16A-17B-18A-19A-20B

\end{document}